% ============================================================================
%  MRes Project Report — DRAFT SKELETON (Introduction section complete)
%
%  IMPORTANT — read before submitting:
%  1. This is an AI-assisted DRAFT. Imperial's regulations require the report to
%     be written in YOUR OWN WORDS. Rewrite every paragraph in your own voice
%     rather than submitting this text as-is, and keep the AI Use Declaration
%     at the end accurate and specific.
%  2. Fill in every <<PLACEHOLDER>> on the title page.
%  3. Figure paths: the \graphicspath below points at a local ./figures/ folder.
%     Copy the PNGs there before compiling (they have been written into your
%     project folder under Project coding/pop_cosmos_notebook/thesis_figures/).
%  4. references.bib is a companion file. One entry (Valiante 2016) is a marked
%     placeholder — verify its DOI/pages on ADS before you cite it.
% ============================================================================

\documentclass[12pt,a4paper]{article}

% ---------- Packages, layout, bibliography (includes.tex) ----------
\input{includes}
\input{notation}

% ---------- Title page metadata ----------
\newcommand{\ReportTitle}{Extending Pop-cosmos to the Far-Infrared:
  A Number-Count Validation}
\newcommand{\degreetype}{Master of Research}
\newcommand{\ProjectCode}{MLBD\_2025\_17}
\newcommand{\StudentName}{Mihir Koka}
\newcommand{\CID}{06067948}
\newcommand{\Supervisor}{David Clements \& Boris Leistedt}
\newcommand{\Assessor}{Alan Heavens}
\newcommand{\WordCount}{<<final word count, excluding layperson's summary>>}
\newcommand{\SubmissionDate}{4 September 2026}

\begin{document}

% ===================== TITLE PAGE =====================
\begin{titlepage}
    \centering

    
\includegraphics[width=0.35\textwidth]{Imperial_College_London_new_logo.png}
    \vspace{1.2cm}

    {\Large Imperial College of Science, Technology and Medicine\par}
    \vspace{0.4cm}
    {\large Department of Physics \par}
    \vspace{1.6cm}

    {\LARGE\bfseries \ReportTitle \par}
    \vspace{1.8cm}

    \begin{tabular}{@{}ll@{}}
        \textbf{Project code:}        & \ProjectCode \\[4pt]
        \textbf{Name:}                & \StudentName \\[4pt]
        \textbf{CID:}                 & \CID \\[4pt]
        \textbf{Supervisor:}          & \Supervisor \\[4pt]
        \textbf{Assessor:}            & \Assessor \\[4pt]
        \textbf{Word count:}          & \WordCount \\[4pt]
        \textbf{Date of submission:}  & \SubmissionDate \\
    \end{tabular}

    \vfill
    {\small MRes in Machine Learning and Big Data in the Physical Sciences\par}
\end{titlepage}

% ===================== LAYPERSON'S SUMMARY =====================
% NOT included in the word count. Limit: one A4 page, < 700 words.
\newpage
\newgeometry{a4paper,margin=2cm}
\section*{\centering Layperson's Summary}
\vspace{-0.3cm}
\begin{center}
{\normalsize\bfseries When Virtual Galaxies Shine Too Brightly}\\[2pt]
{\small Mihir Koka}
\end{center}
\begin{spacing}{1.0}

% --- 1. Opening hook: dust-hidden starlight (~70 words) ---
How do we know a simulated galaxy is realistic when most of its light emerges beyond
human vision? In galaxies where new stars form rapidly, clouds of interstellar
dust absorb the starlight before it reaches our telescopes. The dust warms up and
re-emits that energy as infrared light, at wavelengths several times
longer than visible light. Judging these galaxies from visible light alone means a significant part of this is lost to us

% --- 2. Why the model matters (~80 words) ---
Modern sky surveys record so many galaxies that astronomers need computer-generated
populations to interpret them. Pop-cosmos is one such model; it learns the
properties of nearly half a million galaxies from the COSMOS survey, a deep map of a
patch of sky observed across many wavelengths. It was fitted using visible and
near-infrared light, up to about 4.5 micrometres. My question was whether its
predictions hold much farther into the infrared, at 250 to 500 micrometres. This is
the range where cool interstellar dust, typically only 20 to 40 degrees above absolute
zero, emits most of its energy.

% --- 3. What I did (~100 words) ---
I first checked that the model reproduces observed brightnesses close to the
wavelengths it was trained on. It does, to within a few percent, so any failure at
longer wavelengths would point to the dust physics lacking somewhere.

I then predicted how bright each model galaxy would look to the Herschel Space
Observatory, a telescope launched specifically to observe this cool-dust light, and
counted how many fall into each brightness range per area of sky. I compared those
counts against published results from real Herschel surveys covering several
independent regions. I also tried replacing each galaxy's dust spectrum, the pattern of how its infrared
energy is spread across wavelengths, with warmer alternatives while keeping the total
energy fixed. This isolates the shape of the emission from the amount.

% --- 4. Main findings (~160 words) ---
The model predicts too many bright far-infrared sources. In the 30 to 100 millijansky
range (a unit of brightness used in infrared astronomy), it overestimates by about 1.7 times at 250 and 350 micrometres, rising to about
3.6 times at 500 micrometres. Seven independent datasets agree on this, so it isn't one
survey being unusual.

The cause is cold dust. Model galaxies radiate their peak energy near 136 micrometres,
while observations of galaxies at similar luminosities place the peak closer to 92
micrometres. Colder dust shifts too much energy to the longest wavelengths, overfilling
the 350 and 500 micrometre bins.

A second problem emerged from comparing individual galaxies. The predictions aren't
uniformly wrong; they are scattered. Because bright galaxies are rare, more common faint galaxies scattered upward overfill the bright bins. These two issues are
distinct.

% --- 5. What the modification showed (~60 words) ---
I tested warmer dust spectra while holding each galaxy's total energy fixed.
Intermediate temperatures, around 35 kelvin, matched the published counts considerably
better than the cold baseline or the hottest alternatives. This points toward how pop-cosmos could be improved. The broader lesson seems to be that the model can
match the light it was trained on but has partial failures in the parts it wasn't trained on.
Testing that gap is how we find out what physics is still missing.

\end{spacing}
\restoregeometry

% ===================== ABSTRACT =====================
\newpage
\begin{abstract}
\noindent
The \popcosmos{} framework models the joint distribution of galaxy physical parameters
and photometry, calibrated against optical and near-infrared data from the COSMOS field
\citep{Alsing_2024,Thorp_2025_insights}. Its far-infrared (FIR) luminosity is constrained
by energy balance within the stellar population synthesis model: the total dust luminosity
\LIR{} is fixed by the attenuated starlight, but no galaxy-varying parameter controlling
the \emph{shape} of the FIR spectral energy distribution (SED) is inferred from the
training photometry. FIR predictions are therefore
an untested extrapolation beyond the training range. This work tests them against
observed \emph{Herschel}/SPIRE differential number counts at 250, 350 and \SI{500}{\um},
compiled from nine published count products grouped into three independent survey
families.

Baseline \popcosmos{} overpredicts the counts in all three bands, by a median factor of
$1.24$, $1.72$ and $2.29$ at 250, 350 and \SI{500}{\um} respectively over
\SIrange{10}{30}{\milli\jansky}, rising at brighter fluxes. The excess remains positive
in each of the three independent survey-family benchmarks and exceeds the measured
inter-survey systematic scatter (\dex{0.12}) by a factor of $2.6$--$3.3$. A block
bootstrap over survey families, used as a robustness check rather than a formal
likelihood calculation, gives a median offset of $+0.29$ dex with a 95\% bootstrap
interval of $[+0.16,+0.52]$.

Direct inspection of the model SEDs identifies the cause: the FIR dust bump peaks at a
median rest-frame wavelength of \SI{136}{\um}, with $64\%$ of galaxies peaking beyond
\SI{100}{\um}. Substituting warmer, broader FIR SEDs at fixed \LIR{} --- modified
blackbodies, Casey-like templates \citep{Casey_2012}, and empirical
ALESS/FSPS hybrids --- brings the counts into agreement, though the template families
are mutually degenerate at the \dex{0.03}--\dex{0.09} level and cannot be distinguished
by counts alone.

A per-object comparison against the deblended COSMOS catalogue of \citet{Wang_2024}
appears to contradict the counts, showing the model too \emph{faint} by
\dex{0.16}--\dex{0.24}. The two are reconciled by the model's per-object flux scatter of
$\approx\dex{0.45}$: with count slopes $\alpha = 3.2$--$4.1$, this scatter promotes
intrinsically faint galaxies into bright flux bins, inflating bright counts by
\dex{0.55}--\dex{0.86} while typical galaxies remain underluminous. The two effects
nearly cancel, reproducing the observed net excess. This scatter is unchanged by every
template variant tested, identifying two independent deficiencies --- an SED shape
problem that a warmer dust prescription resolves, and a flux-scatter problem that it
does not.
\end{abstract}

% ===================== INTRODUCTION =====================
\section{Introduction}
\label{sec:intro}

\subsection{Why far-infrared light matters}
\label{sec:intro:fir}

Roughly half of all starlight ever emitted has been absorbed by interstellar dust and
re-radiated in the infrared \citep{Dole_2006}. The
obscured fraction rises with star-formation rate (SFR), so the galaxies
contributing most to cosmic star-formation history are the ones whose ultraviolet and
optical light is most heavily attenuated. Any dataset of star formation built on
short-wavelength data alone is insufficient and incomplete in a way that depends on the
population being studied.

Dust absorbs at short wavelengths and re-emits as thermal continuum peaking, for
typical star-forming galaxies, between roughly 30 and \SI{200}{\um} in the rest frame.
A galaxy's spectral energy distribution (SED), the way its energy is spread across
wavelengths, therefore has a peak in the far-infrared that controls both how much dust
is present and how hot it is. The
\emph{Herschel} Space Observatory \citep{Pilbratt_2010} provided the most complete
extragalactic far-infrared survey coverage to date through its SPIRE instrument
\citep{Griffin_2010}, which operated at 250, 350 and \SI{500}{\um}. These
wavelengths sit close to the dust emission peak for galaxies at $z \sim 1$--$3$, so
SPIRE data constrain the total obscured energy output of the population directly.

\subsection{From individual galaxies to population models}
\label{sec:intro:population}

Modern sky surveys record so many galaxies that per-object SED fitting becomes both
computationally expensive and statistically difficult. Individual posteriors are
prior-dominated for faint sources, and combining them into population-level statements
can't be done simply. Population-level generative modelling offers an alternative. Rather than
fitting each galaxy separately, one models the joint distribution of physical parameters
across the population and forward-models the photometry, allowing selection effects,
noise and calibration to be treated self-consistently \citep{Alsing_2024}.

Surveys such as COSMOS \citep{Scoville_2007} combine many wavelength bands over the
same area of sky. Conventional SED-fitting tools like \textsc{fsps}
\citep{Conroy_2009} and \textsc{Prospector} \citep{Leja_2017} infer quantities such as
stellar mass, SFR, dust attenuation and redshift from this photometry one galaxy at a
time. \popcosmos{} \citep{Alsing_2024,Thorp_2024} takes a different approach. It
describes the joint distribution of galaxy redshift and physical parameters using a
score-based diffusion model, so that correlated properties (stellar mass,
star-formation history, metallicity, dust attenuation and AGN contribution) are put
together rather than taken from unrelated one-dimensional distributions. A 16-parameter
stellar population synthesis model then converts each parameter vector into a
rest-frame SED and predicted photometry. A survey data model adds calibration effects,
photometric uncertainties and the same selection function applied to the COSMOS2020
catalogue. Synthetic and observed 26-band catalogues are compared using an
optimal-transport distance, allowing the population and nuisance models to be
calibrated together. The mathematical details of diffusion training and optimal
transport are given in the original papers \citep{Alsing_2024,Thorp_2024}.

This makes \popcosmos{} useful for survey simulation and
forecasting, but it also makes validation beyond the calibration regime important. A plausible
catalogue in the fitted bands is not automatically correct at unseen wavelengths.

\subsection{The far-infrared extrapolation problem}
\label{sec:intro:extrapolation}

The photometry used to constrain \popcosmos{} spans roughly \SIrange{0.3}{5}{\um},
optical to near-infrared. The model receives no far-infrared information. It still makes
far-infrared predictions, because dust emission in \textsc{fsps} is tied to attenuation
by \emph{energy balance}: the total luminosity absorbed by dust in the ultraviolet and
optical is re-emitted in the infrared, conserving energy. This fixes the integrated dust
luminosity \LIR{} for every galaxy.

The key point for this project is what energy balance does not fix. It constrains
the total infrared output but not its shape. The wavelength at which the dust emission
peaks depends on the dust temperature, and the optical/near-infrared training
data is not enough to perfectly model this. The far-infrared predictions of \popcosmos{} are therefore an
extrapolation, one that has not been previously tested against far-infrared observations.
A disagreement here
is useful to study because it probes physics that the calibration data couldn't
constrain.

\subsection{Why use observed number counts}
\label{sec:intro:counts}

Validating a population model in a new wavelength regime requires having some comparison
metric. An obvious approach is to compare derived physical quantities (\LIR{} or
obscured SFR) against published catalogue values. In practice this is awkward becase catalogue
\LIR{} values are themselves outputs of SED fits with their own dust templates and
initial mass function conventions, so comparing physical quantities tests the union of
two modelling chains rather than the model itself.

Differential number counts avoid most of this. The count $\mathrm{d}N/\mathrm{d}S$, the
surface density of sources per unit flux density (brightness measured in one observing
band), is measured in observed flux space. The model prediction requires only
redshifting each galaxy's SED and evaluating its flux at the SPIRE band-centre
wavelength. No dust-template inversion, initial mass function conversion or timescale
reconciliation is involved. Counts are a long-established diagnostic with a substantial
published literature at SPIRE wavelengths \citep{Clements_2010, Oliver_2010,
Glenn_2010, B_thermin_2010}, providing multiple independent measurements against which a
model can be tested.

Two properties of counts need care. First, they're conventionally reported in
Euclidean-normalised form, $S^{2.5}\,\mathrm{d}N/\mathrm{d}S$, which flattens the
steep intrinsic slope and makes the differences visible. This normalisation is used
throughout. Second, differential counts are preferred over cumulative counts $N(>S)$
because cumulative bins re-use sources at every threshold and therefore have strongly
correlated errors. Counts aren't free of systematics either as they require completeness
and flux-boosting corrections, and in confusion-dominated SPIRE maps those corrections
are substantial \citep{Nguyen_2010, Roseboom_2010}.
Section~\ref{sec:data} distinguishes published corrected counts, used as
the benchmark, from raw deblended catalogue counts, used only as a diagnostic.

\subsection{Aim, questions and scope}
\label{sec:intro:aims}

This project was originally framed as a multi-wavelength validation of \popcosmos{}
spanning far-infrared, radio and X-ray. However early work showed that each regime
required resolving its own systematics, and that the far-infrared alone contained
enough substance for a coherent project. The scope was narrowed to far-infrared on the
recommendation that validation focus on directly observed quantities. The objectives
are:

\begin{enumerate}
  \item Assemble a benchmark of published SPIRE differential number counts at 250, 350
        and \SI{500}{\um} from independent sky fields, in consistent units, and
        establish which sources can be treated as statistically independent
        (Section~\ref{sec:data}).
  \item Predict \popcosmos{} number counts by propagating each model galaxy's SED to
        observed flux, and quantify agreement with the benchmark using a metric whose
        error budget reflects the systematic scatter between surveys
        (Section~\ref{sec:methods}).
  \item Determine whether any disagreement can be attributed to the far-infrared SED
        shape, by substituting alternative dust prescriptions at fixed \LIR{} and
        measuring the resulting change in the counts (Section~\ref{sec:methods}).
  \item Test robustness against correlated residuals, the assumed error floor,
        small-number statistics at the bright end, and known systematics in the
        comparison catalogues (Section~\ref{sec:results}).
  \item Determine why the model fails by isolating the discrepancy in redshift
        and considering the count-level result with a per-object flux comparison
        (Section~\ref{sec:results}).
\end{enumerate}

% ===================== DATA SECTION =====================
\section{Data}
\label{sec:data}

\subsection{\popcosmos{} catalogue and generated SEDs}
\label{sec:data:popcosmos}

The work in this report uses the infrared-selected \popcosmos{} catalogue of 429,669
galaxies drawn from the COSMOS2020 photometric survey
\citep{Thorp_2025_insights}. Each galaxy has
median posterior estimates of 16 stellar population synthesis parameters, including
redshift ($0 \leq z \leq 6$, median $z = 1.28$), stellar mass, star-formation history,
dust attenuation, gas-phase metallicity and AGN bolometric luminosity fraction
$f_\mathrm{AGN}$. The model was fitted against 26-band photometry spanning roughly
\SIrange{0.3}{4.5}{\um} (observed frame), from the $u$ band through IRAC channels~1
and~2 \citep{Thorp_2025_insights}.

For every galaxy, a full rest-frame SED is generated using the \textsc{fsps} stellar
population synthesis code and its stored median posterior parameters. The SED data include
the attenuated stellar$+$dust spectrum sampled at 5,994 rest-frame wavelength points from
\SI{0.009}{\um} to \SI{10\,000}{\um} and the integrated dust luminosity
\LIR{} ($8$--$\SI{1000}{\um}$). Wavelengths in the original files are in \AA ngstr\"om
and were converted to micrometres throughout the entire project. All SEDs used in this work are derived from median posterior parameters instead of full posteriors. This is sufficient
for an initial population-level validation but breaks partially with some
parameters, particularly $f_\mathrm{AGN}$, whose prior is bimodal
\citep{Thorp_2025_insights}. The consequences of this choice are discussed further in
Section~\ref{sec:discussion}.

When converting model galaxies into predicted number counts, the sky area of
the field must be specified. The Farmer COSMOS2020 catalogue flag
\texttt{FLAG\_COMBINED} defines the area which covers
\SI{1.278}{\deg\squared}. This value was used for all model count predictions.

\subsection{Published SPIRE differential counts}
\label{sec:data:counts}

The benchmark of observed counts consists of published, corrected differential
number counts at 250, 350 and \SI{500}{\um} compiled from six studies. These are
summarised in Table~\ref{tab:counts} and shown in Figure~\ref{fig:compiled_counts}.

\begin{table}[H]
\centering
\caption{Published SPIRE count datasets used. The primary surveys
(H-ATLAS, HerMES, Dark Field) cover independent sky areas and provide the main
quantitative benchmark. The remaining tables overlap these fields and serve as
sanity checks on the methods.}
\label{tab:counts}
\footnotesize
\setlength{\tabcolsep}{4pt}
\begin{tabular}{@{}llll@{}}
\toprule
Reference & Field & Method & Flux (mJy)  \\
\midrule
\citet{Valiante_2016} & H-ATLAS GAMA   & matched filter    & 25--300   \\
\citet{Oliver_2010}   & HerMES         & source extraction & 24--374   \\
\citet{Pearson_2025}  & SPIRE Dark Fld & XID deblending    & 11--100   \\
\midrule
\citet{Clements_2010} & H-ATLAS SDP    & source extraction & 34--800  \\
\citet{Glenn_2010}    & HerMES         & P(D) fluctuation  & 2--1000  \\
\citet{Varnish_2025}  & SPIRE Dark Fld & P(D) fluctuation  & 0.02--1000 \\
\bottomrule
\end{tabular}
\end{table}

Although the six papers provide 235 transcribed count rows across all three bands, they
do not all serve as independent measurements. \citet{Valiante_2016} covers the wide H-ATLAS GAMA
fields and provides the strongest bright-end constraint. \citet{Oliver_2010} covers the
separate HerMES multi-tier fields and is most useful in the middle flux range.
\citet{Pearson_2025} covers the SPIRE Dark Field near the North Ecliptic Pole and
extends the counts to the faintest fluxes. The remaining three studies overlap
these fields. \citet{Clements_2010} uses an earlier H-ATLAS release of the same sky as
Valiante. \citet{Glenn_2010} applies a P(D) fluctuation analysis to HerMES maps, and
\citet{Varnish_2025} applies P(D) to the same Dark Field maps used by Pearson.

% FIGURE: compiled published counts
\begin{figure}[H]
\centering
\makebox[\textwidth][c]{\includegraphics[width=1.2\textwidth]{Figures/compiled_counts_thesis.png}}
\caption{Published Euclidean-normalised differential counts at 250, 350 and
\SI{500}{\um} from the six studies listed in Table~\ref{tab:counts}. The main
 comparison in this work uses the 10--\SI{100}{mJy} range (shaded),
where all the surveys agree closely and the counts are well constrained.
Bins above \SI{100}{mJy} are shown but treated with caution because of small sample sizes due to
the limited COSMOS area and the increasing effect of gravitationally lensed
sources.}
\label{fig:compiled_counts}
\end{figure}

The source papers use a mixture of steradian- and square-degree-based units. For this
analysis all counts were converted to the common Euclidean-normalised form
$S^{2.5}\,\mathrm{d}N/\mathrm{d}S$, with units of
$\mathrm{Jy}^{1.5}\,\mathrm{deg}^{-2}$. P(D) counts from \citet{Glenn_2010} and
\citet{Varnish_2025} are statistical constraints derived from pixel-brightness
fluctuations below the individual-source detection limit. They are not catalogued
galaxies and thus are not used for per-object comparisons.

\subsection{Wang and Jin COSMOS catalogues}
\label{sec:data:catalogues}

Two deblended far-infrared catalogues of the COSMOS field are used for per-object flux
comparisons. \citet{Wang_2024} provides a probabilistic, prior-based deblended catalogue
containing 131,178 rows, of which 128,387 are associated with positive COSMOS2020
identifiers; the remaining 2,791 are radio-only priors. SPIRE flux densities were
extracted using the XID+ framework \citep{Hurley_2017}. \citet{Jin_2018} provides a
super-deblended catalogue of 195,107 sources with a similar extraction process.

These catalogues were used for a different purpose than the published counts. They provide
individual galaxy flux measurements and uncertainties, which are needed for the
per-object scatter analysis in Section~\ref{sec:results}. They are not
independent sky fields, and their raw entries can't be compared directly to the completeness-corrected count tables from above. The purpose of having both is for their
agreement with each other. Wang and Jin fluxes for matched sources agree to
0.12--\SI{0.15}{dex}, which is much better than either catalogue's agreement with
\popcosmos{} (0.44--\SI{0.54}{dex}). This allows the large per-object scatter to be isolated to
the model rather than to one deblending pipeline.

\subsection{Data-quality and selection limits}
\label{sec:data:quality}

For per-object comparisons against the Wang catalogue, a signal-to-noise ratio (SNR)
threshold of $\geq 5$ is applied to the Wang flux measurements. This conservative choice
is motivated by \citet{Clements_1999}, who note that flux-threshold bias and weak-source reliability are
more easily controlled in catalogues detected at $\geq 5\sigma$.
The threshold removes fewer than 3\% of sources at \SI{20}{mJy} but a larger amount
at fainter fluxes and longer wavelengths.

The published count tables used as the primary benchmark were already corrected for
completeness and flux boosting by their authors. The main count claims in this work focus on
the approximately 10--\SI{100}{mJy} range, where all three primary survey families
overlap and the corrections are best characterised. Bins above approximately
\SI{100}{mJy} are shown but should be considered with caution, because the COSMOS area
(\SI{1.278}{\deg\squared}) provides few bright sources and the greater effect of lensed
objects makes comparisons at these longer wavelengths more complicated.

\section{Methods}
\label{sec:methods}

\subsection{From a rest-frame SED to an observed flux}
\label{sec:methods:flux}

Each galaxy in the \popcosmos{} catalogue has a rest-frame attenuated SED
$L_\nu^\mathrm{rest}(\lambda)$ in units of $L_\odot\,\mathrm{Hz}^{-1}$, stored on a
grid of 5,994 rest-frame wavelengths. To predict the observed SPIRE flux density at an observed wavelength $\lambda_\mathrm{obs}$, the
rest-frame wavelength probed is
\begin{equation}
  \lambda_\mathrm{rest} = \frac{\lambda_\mathrm{obs}}{1+z}\,,
  \label{eq:restlam}
\end{equation}
and the observed flux density is
\begin{equation}
  S_\nu = \frac{(1+z)\,L_\nu^\mathrm{rest}(\lambda_\mathrm{rest})}
               {4\pi\,d_L^2(z)}\,,
  \label{eq:flux}
\end{equation}
where $d_L(z)$ is the luminosity distance. Luminosity distances are computed using the
WMAP9 cosmology ($H_0 = 69.32\;\mathrm{km\,s^{-1}\,Mpc^{-1}}$,
$\Omega_m = 0.2865$), adopted consistently throughout this conversion. The factor
$(1+z)$ converts between rest-frame and observed frequency
intervals (the $K$-correction). $L_\nu^\mathrm{rest}$ is evaluated at
$\lambda_\mathrm{rest}$ by linear interpolation on the stored wavelength grid, and the
result is converted to millijanskys.

This calculation samples the SED at each nominal SPIRE wavelength rather than integrating
through the full filter transmission curve. It is a computational approximation applied
consistently to every template. Full bandpass integration could shift individual fluxes
and should be included in future precision or per-object analyses.

As a reproducibility check, the integrated $8$--$\SI{1000}{\um}$ luminosity computed
from the stored SED was compared against the stored \LIR{} value. The median ratio
across a random sample of 500 galaxies is $1.00002$, confirming internal consistency of
the SED products.

Later in this work, a median population SED is used; it is constructed by placing
all SEDs in the selected population on the same rest-frame wavelength grid, normalising
each to its \LIR{}, and then taking the median and 16th--84th percentile
interval. This is different from just plotting some individual galaxy whose \LIR{} happens to
fall at the population median.

\subsection{Constructing differential number counts}
\label{sec:methods:counts}

For each SPIRE band the predicted fluxes are binned into logarithmically spaced
intervals. The differential count in each bin is
\begin{equation}
  \frac{\mathrm{d}N}{\mathrm{d}S} = \frac{N_\mathrm{bin}}{\Delta S \;\Omega}\,,
\end{equation}
where $N_\mathrm{bin}$ is the number of model galaxies in the bin, $\Delta S$ is the
bin width in janskys, and $\Omega = \SI{1.278}{\deg\squared}$ is the effective catalogue
area. Counts are reported in Euclidean-normalised form,
$S^{2.5}\,\mathrm{d}N/\mathrm{d}S$, with units of
$\mathrm{Jy}^{1.5}\,\mathrm{deg}^{-2}$. As discussed previously in the paper this is to make the disagreements between the model and data easier to spot across the flux range. It is just a rescaling applied to both the model and published curves.

\subsection{Alternative far-infrared SED families}
\label{sec:methods:templates}

The main question here is whether the count disagreement improves when using a warmer far-infrared
shape while the learned galaxy population and \LIR{} are held fixed. To answer this, four
alternative dust prescriptions are substituted in place of the FSPS baseline. In all tests, each galaxy's total infrared luminosity is preserved, so the
comparison changes only where in wavelength the energy appears, not how much is emitted.
These alternatives were tested as a purely exploratory exercise as they do show the direction of the needed correction but do not directly identify the underlying physics at play. 

The FSPS baseline is the existing \popcosmos{}/FSPS dust emission shape, used
without modification. The ALESS template is the composite rest-frame
SED of luminous submillimetre galaxies from \citet{daCunha_2015}, renormalised to each
galaxy's \LIR{}. It serves as a warm-shape bracket rather than a universal galaxy
model. FSPS/ALESS hybrids were created by combining the two luminosity-density shapes on a
common rest-frame wavelength grid at 25\%, 50\% and 75\% ALESS contribution, each
renormalised to the original \LIR{}. The modified-blackbody (MBB) grid uses optically
thin, single-temperature shapes, $L_\nu \propto \nu^\beta B_\nu(T)$, with dust
temperatures from 20 to \SI{50}{K} in \SI{5}{K} steps. The emissivity index is fixed at
the assumed value $\beta=1.8$, rather than inferred from the data. Each template is
renormalised over 8--\SI{1000}{\um} to preserve the galaxy's original \LIR{}.
Finally, a Casey-inspired diagnostic replaces the Wien side of the MBB with a continuous
mid-infrared power law, joined where the logarithmic slopes match
\citep{Casey_2012}. This captures the greybody-plus-power-law idea but is not the exact
analytic Casey fitting function. It is treated as a diagnostic parameterisation, not a
complete solution that explains the physics. Figure~\ref{fig:mbb_shapes} shows the modified
blackbody shapes normalised to equal \LIR{}, illustrating how the same total energy
is redistributed across wavelength as dust temperature increases.

% FIGURE: MBB template shapes
\begin{figure}[H]
\centering
\makebox[\textwidth][c]{\includegraphics[width=1.15\textwidth]{Figures/popcosmos_mbb_temperature_grid_shapes.png}}
\caption{Modified blackbody dust spectra at temperatures from 20 to \SI{50}{K}
($\beta = 1.8$), each normalised to the same integrated \LIR{}. Warmer dust shifts
the emission peak to shorter wavelengths and reduces the flux at 350 and
\SI{500}{\um} relative to the total. The FSPS baseline peak falls near the 20--25~K
curves.}
\label{fig:mbb_shapes}
\end{figure}

\begin{figure}[H]
\centering
\makebox[\textwidth][c]{\includegraphics[width=1.15\textwidth]{Figures/popcosmos_mbb_temperature_grid_counts.png}}
\caption{Predicted differential counts for the modified blackbody temperature grid,
compared against published observations. Each curve uses the same \LIR{} per galaxy;
only the dust temperature changes. Intermediate temperatures around 30--\SI{35}{K}
bring the predicted counts closest to the observations, while colder and hotter
extremes both produce worse agreement.}
\label{fig:mbb_counts}
\end{figure}

% FIGURE: analysis pipeline schematic
\begin{figure}[H]
\centering
\makebox[\textwidth][c]{\includegraphics[width=1.15\textwidth]{Figures/popcosmos_fir_validation_pipeline.png}}
\caption{Main work pipeline. \popcosmos{} parameters and redshifts are converted to
rest-frame SEDs via FSPS, then to observed SPIRE fluxes. Differential number counts are compared against published
survey results (left), and per-object fluxes are compared against the Wang and Jin
deblended catalogues (right).}
\label{fig:pipeline}
\end{figure}

\subsection{Count evaluator and error treatment}
\label{sec:methods:evaluator}

Model and published counts are compared using a chi-square evaluator operating in
logarithmic count space. For each published flux bin $i$ with observed
Euclidean-normalised count $C_i^\mathrm{obs}$ and uncertainty $\sigma_i$, the model
count $C_i^\mathrm{mod}$ is obtained by interpolating the model curve to the same flux
value. The log-space residual is
\begin{equation}
  \Delta_i = \log_{10} C_i^\mathrm{mod} - \log_{10} C_i^\mathrm{obs}\,,
\end{equation}
A positive $\Delta_i$ means the model overpredicts. Published uncertainties vary widely in how they are estimated, and some studies quote very small formal errors that do not reflect the true systematic scatter between surveys. An error floor of \SI{0.12}{dex} is therefore imposed, corresponding to the
measured inter-survey scatter between independent count families at the same flux. The
primary statistic to evaluate is the median $\Delta_i$ across bins.
A reduced $\chi^2$ is also reported as a rough ranking
measure but is not treated as a formal likelihood, because adjacent flux bins within
the same survey share calibration and completeness corrections and are therefore
correlated. \\

Each bin is one flux interval from one published table, not an individual galaxy.
Because the six studies observe overlapping sky areas
(Section~\ref{sec:data:counts}), treating all of them as independent would overcount
the evidence. For this reason, the chi-square is computed on two sets. The full set uses 174
bins from all six studies. The clean set uses only the 74 bins from Valiante, Oliver
and Pearson, which cover three separate patches of sky and provide the
main test. The model ranking is the same on both sets. As a further check, the
comparison is repeated with a larger error floor and with a block bootstrap that
resamples across the three independent survey families.

\subsection{Object-level residuals and the scatter test}
\label{sec:methods:scatter}

For galaxies matched between the \popcosmos{} predictions and the Wang or Jin
catalogues, per-object residuals are defined as
\begin{equation}
  r = \log_{10}\!\left(\frac{F_\mathrm{model}}{F_\mathrm{observed}}\right).
\end{equation}
Only sources with observed SNR~$\geq 5$ in the relevant band are kept. The median
$r$ measures systematic offset, and the 16--84th percentile half-width measures
scatter. Three comparisons are computed (model vs Wang, model vs Jin, Wang vs
Jin) to separate model scatter from deblending scatter.

The scatter has a direct consequence for the counts. With a decreasing source
distribution, symmetric per-object scatter is not consistent in its effect on number
counts, because more faint sources are scattered upward than bright sources are
scattered downward. To actually quantify this, a test that selects model galaxies
predicted above a flux cut was conducted (e.g.\ \SI{20}{mJy}). This asks how many are also observed
above that cut. The ratio shows how many model-bright sources are actually bright
versus those just promoted by scatter. This effect is commonly described as
Eddington bias or flux boosting. Related flux-limited selection effects are also called
Malmquist bias in parts of the source-count literature \citep{Clements_1999}.

\subsection{Physical diagnostics}
\label{sec:methods:diagnostics}
Two more diagnostics are looked at beyond the overall count mismatch and ask where the model goes wrong. The first is related to how real galaxies follow a pattern in which brighter galaxies have warmer dust that peaks at shorter wavelengths. Drew and Casey (2022) calibrated this relationship and showed it holds without redshift evolution out to $z \sim 2$. For each low-AGN ($f_\mathrm{AGN} < 0.1$) pop-cosmos galaxy, the rest-frame far-infrared peak wavelength is measured and plotted against infrared luminosity. If the model worked perfectly, it would place the right amount of energy at SPIRE wavelengths as a function of luminosity. If it does not, the comparison would show whether the dust is too cold, too uniform, or both.

The second test asks which galaxies are responsible for the count excess. Model galaxies are split into redshift bins, and one bin is removed at a time. The bin whose removal reduces the excess by the largest amount is where the problem lies most. This matters because the same observed SPIRE band samples different rest-frame wavelengths at different redshifts, so the cold-dust problem wouldn't necessarily affect all epochs equally.

\section{Results}
\label{sec:results}

\subsection{Control test where the model is constrained}
\label{sec:results:control}

Before testing the far-infrared predictions, it is important to confirm that
\popcosmos{} performs well where the input photometry constrains it. Predicted and
observed IRAC magnitudes were compared for all galaxies with coverage in each channel.
At Ch1 (\SI{3.6}{\um}), 423,272 galaxies yield a median residual of $+0.020$ mag with
a median absolute deviation (MAD) of 0.084 mag. At Ch2 (\SI{4.5}{\um}), 414,272
galaxies yield $-0.012$ mag (MAD 0.041 mag). These residuals are small enough to
confirm that the catalogue, the SED products and the comparison pipeline are working
correctly. Ch3 and Ch4 have lower coverage and are weaker controls.

% FIGURE: IRAC control test
\begin{figure}[H]
\centering
\makebox[\textwidth][c]{\includegraphics[width=1.15\textwidth]{Figures/popcosmos_irac_redshift_histograms.png}}
\caption{Predicted vs observed IRAC Ch1 and Ch2 magnitudes as a function of redshift.
The model reproduces the near-infrared photometry to within a few hundredths of a
magnitude across the full redshift range. This confirms that the far-infrared mismatch
reported below is wavelength-dependent, not a universal catalogue failure.}
\label{fig:irac_control}
\end{figure}

The key implication is that any disagreement at 250--\SI{500}{\um} points to the
far-infrared dust physics specifically, not to a systematic problem with the population
model or the flux-prediction pipeline.

\subsection{Published far-infrared number counts}
\label{sec:results:counts}

Figure~\ref{fig:count_overlay} shows the headline result. Baseline \popcosmos{}/FSPS
differential counts are plotted against the published observations from all six studies
at 250, 350 and \SI{500}{\um}.

% FIGURE: headline count overlay
\begin{figure}[H]
\centering
\makebox[\textwidth][c]{\includegraphics[width=1.15\textwidth]{Figures/popcosmos_count_overlay_thesis.png}}
\caption{Euclidean-normalised differential counts at 250, 350 and \SI{500}{\um}.
Coloured points are published observations (Table~\ref{tab:counts}); the solid line is
the baseline \popcosmos{}/FSPS prediction. The model overpredicts in all three bands,
with the excess growing toward longer wavelengths. Dashed and dotted lines show
alternative dust prescriptions at fixed \LIR{} (Section~\ref{sec:results:templates}).}
\label{fig:count_overlay}
\end{figure}

The baseline model overpredicts the counts in all three bands. In the well-constrained
30--\SI{100}{mJy} range, the excess is approximately $1.7\times$ at both 250 and
\SI{350}{\um}, rising to approximately $3.6\times$ at \SI{500}{\um}. Expressed as
median log residuals across the full independent-source scorecard, the offsets are
$+0.25$, $+0.41$ and $+\SI{0.55}{dex}$ at 250, 350 and \SI{500}{\um} respectively.
These two summaries use different flux ranges and pooling methods; the 30--100 mJy
factors come from the regime where all three primary survey families overlap, while the
median residuals span the full 10--300 mJy range. Both measures agree that the excess
is positive in every independent count source, grows with wavelength, and exceeds the
measured inter-survey scatter of \SI{0.12}{dex} by a factor of 2.6--3.3.

\subsection{The baseline SED is too cold and too uniform}
\label{sec:results:cold}

The count overprediction at long wavelengths implies that the model places too much
energy beyond the SPIRE peak. To test this independently of the template experiment, the
rest-frame far-infrared peak wavelength $\lambda_\mathrm{peak}$ of each low-AGN
($f_\mathrm{AGN} < 0.1$) model galaxy was measured and compared against infrared
luminosity.

% FIGURE: Drew & Casey peak relation
\begin{figure}[H]
\centering
\makebox[\textwidth][c]{\includegraphics[width=1.15\textwidth]{Figures/drew_casey_peak_relation.png}}
\caption{Rest-frame far-infrared peak wavelength vs infrared luminosity for low-AGN
\popcosmos{} galaxies (grey density), compared with the observed relation from
\citet{Drew_2022} (red line). The observed relation has a negative slope
($\eta = -0.09 \pm 0.01$), meaning more luminous galaxies have warmer, shorter-peaking
dust. The model slope is essentially flat ($\eta = +0.005$), and the peak is offset to
longer wavelengths (approximately \SI{126}{\um} vs \SI{92}{\um} at
$L_\mathrm{IR} = 10^{12}\,L_\odot$).}
\label{fig:drew_casey}
\end{figure}

\citet{Drew_2022} calibrated an empirical relation with no redshift evolution out to
$z \sim 2$, finding
$\lambda_\mathrm{peak} \propto L_\mathrm{IR}^{-0.09}$, with
$\lambda_\mathrm{peak} \approx \SI{92}{\um}$ at
$L_\mathrm{IR} = 10^{12}\,L_\odot$. The \popcosmos{} model shows two failures against
this benchmark. First, a normalisation offset. Model galaxies peak near \SI{126}{\um}
at the same luminosity, corresponding to colder dust. Second, and more significantly,
the model slope is essentially flat ($\eta = +0.005$), meaning the dust peak does not
shift with luminosity at all. The model also has an unusually narrow range of peak
wavelengths. The missing luminosity dependence means the bright, luminous population
remains too cold at exactly the wavelengths where the SPIRE counts are most sensitive.

The primary peak-relation test uses low-AGN objects because the population SED
diagnostic revealed a separate hot-peaking branch controlled by the model
$f_\mathrm{AGN}$ parameter. Approximately 67\% of objects with $f_\mathrm{AGN} > 0.3$
peak below \SI{40}{\um}, and 94\% of all hot-peaking objects lie in that high-AGN
group. This supports interpreting the unusual 20--\SI{40}{\um} peaks as model AGN/torus
emission rather than ordinary cold dust. However, $f_\mathrm{AGN}$ is a model
luminosity-ratio parameter, not an external AGN detection, and its bimodal prior makes
posterior-median interpretation delicate \citep{Thorp_2025_insights}. The low-AGN cut
prevents this hot branch from driving the claim that the star-forming far-infrared peak
is too cold. The detailed AGN diagnostic is presented in Appendix~B.

\subsection{Controlled template experiments}
\label{sec:results:templates}

To confirm that the count excess responds to a warmer far-infrared shape, the
alternative dust prescriptions described in Section~\ref{sec:methods:templates} were
substituted at fixed \LIR{}. The resulting count curves are included in
Figure~\ref{fig:count_overlay} as dashed/dotted lines.

Warmer templates reduce the excess in the expected direction. The best-scoring
alternatives in the unified 174-point comparison are Casey $T = \SI{30}{K}$,
$\alpha = 2.5$ (median residual $-\SI{0.06}{dex}$) and Casey $T = \SI{30}{K}$,
$\alpha = 3.0$ ($+\SI{0.03}{dex}$). The MBB at \SI{35}{K} ($-\SI{0.12}{dex}$) and the
50\% ALESS/FSPS hybrid ($-\SI{0.08}{dex}$) also perform well. All four are consistent
with zero within the inter-survey scatter floor. In the cleaner 74-point three-family
evaluation, the ranking is similar, with Casey $T = \SI{30}{K}$, $\alpha = 3.0$ scoring
best (median $+\SI{0.01}{dex}$).

The improvement is robust but not uniquely diagnostic. Several template families score
similarly within inter-survey scatter, so the counts identify a needed warmer and more
luminosity-dependent direction but do not uniquely select a dust prescription. The FSPS
baseline remains positive in the survey-family block bootstrap (median
$+\SI{0.29}{dex}$; 95\% bootstrap interval $[+0.16,+0.52]$), while the best-scoring
alternatives straddle zero. This interval is a robustness summary rather than a formal
95\% exclusion because only three independent survey families are available.

\subsection{Per-object scatter is a second problem}
\label{sec:results:scatter}

A separate finding emerges from comparing individual galaxies rather than population
counts.

% FIGURE: scatter mechanism
\begin{figure}[H]
\centering
\makebox[\textwidth][c]{\includegraphics[width=1.15\textwidth]{Figures/scatter_mechanism_explained.png}}
\caption{Left: predicted vs observed 250~\SI{}{\um} flux for matched galaxies, showing
the large per-object scatter. Right: among 422 model galaxies predicted above
\SI{20}{mJy}, only 21 are also observed above that threshold; 401 are promoted into
the bright regime by scatter alone. In a steeply falling source distribution, this
asymmetric promotion inflates the bright-end counts even when the median prediction is
too faint.}
\label{fig:scatter}
\end{figure}

Matched against the Wang catalogue, model galaxies show per-object scatter of
0.44--\SI{0.54}{dex} (16--84th percentile half-width) at 250, 350 and \SI{500}{\um}.
This is far larger than the 0.12--\SI{0.15}{dex} agreement between the Wang and Jin
catalogues for the same sources, confirming that the scatter originates in the model
rather than the deblending. The scatter is stable under SNR cuts at 3, 5, 7 and
$10\sigma$, so observational noise near the detection limit does not drive it.

The scatter has a direct consequence for the counts, which resolves an apparent paradox.
Matched detections have a median model flux that is too \emph{faint} (by
approximately $-\SI{0.16}{dex}$ at 250~$\upmu$m), yet the full model catalogue produces
too many bright sources. Both statements are true simultaneously. In a steeply falling
source distribution, symmetric scatter promotes more faint galaxies upward across a
bright threshold than it moves bright galaxies downward. Among 422 model galaxies
predicted above \SI{20}{mJy} at 250~$\upmu$m, only 21 are also observed above that cut;
401 are promoted by scatter alone. A decomposition of the count residuals into
contributions from the mean offset and from the scatter
(Table~\ref{tab:paradox}) shows that at 250~$\upmu$m the scatter effect
($+\SI{0.55}{dex}$) fully accounts for the measured excess ($+\SI{0.19}{dex}$), with
the mean offset partially cancelling it. At 350 and 500~$\upmu$m the scatter accounts
for roughly half to two-thirds of the excess, and the cold-dust mean offset contributes
the remainder.

\begin{table}[H]
\centering
\caption{Decomposition of the count excess into contributions from the per-object mean
offset and from scatter, at each SPIRE wavelength.}
\label{tab:paradox}
\footnotesize
\begin{tabular}{@{}rrrrrr@{}}
\toprule
Band & Mean offset & Count effect & Scatter effect & Predicted net & Measured net \\
($\upmu$m) & (dex) & of mean (dex) & (dex) & (dex) & (dex) \\
\midrule
250 & $-0.163$ & $-0.361$ & $+0.549$ & $+0.19$ & $+0.19$ \\
350 & $-0.235$ & $-0.621$ & $+0.731$ & $+0.11$ & $+0.31$ \\
500 & $-0.183$ & $-0.575$ & $+0.855$ & $+0.28$ & $+0.43$ \\
\bottomrule
\end{tabular}
\end{table}

A warmer mean template can move the counts, but it does not by itself remove the
object-to-object scatter. The two problems are distinct.

\subsection{Where the excess occurs in redshift}
\label{sec:results:redshift}

% FIGURE: z-resolved diagnosis
\begin{figure}[H]
\centering
\makebox[\textwidth][c]{\includegraphics[width=1.15\textwidth]{Figures/z_resolved_fir_diagnosis.png}}
\caption{Change in median count residual when each redshift bin is excluded from the
model prediction. Larger improvement (taller bar) means that redshift bin contributes
more to the excess. At 250~$\upmu$m the excess is concentrated at $z < 1$; at
500~$\upmu$m it shifts to $z \approx 1.5$--$2.5$.}
\label{fig:z_resolved}
\end{figure}

The excess does not come from one universal redshift range. At 250~$\upmu$m, the
largest improvement comes from excluding $z < 1$ (improvement of
\SI{0.28}{dex}). At 500~$\upmu$m, the excess shifts to $z \approx 1.5$--$2.5$
(improvement of \SI{0.22}{dex}). This band-dependent localisation is expected, because
the same observed wavelength samples different rest-frame wavelengths as redshift
changes. \citet{B_thermin_2012} found from a stacking analysis that the characteristic
contribution to SPIRE counts shifts to higher redshift at longer wavelengths, which is
consistent with this pattern. A direct numerical comparison against their
redshift-sliced counts was not performed, as those tabulated values were not imported
into the evaluator; this remains a valuable extension.

\subsection{Observational side result: the Valiante bright-end uptick}
\label{sec:results:valiante}

While checking the observational inputs, an unexpected feature was found in the
\citet{Valiante_2016} H-ATLAS counts. In the brightest bins (above approximately
\SI{200}{mJy}), the Euclidean-normalised counts rise rather than continuing to fall.
The final \SI{300}{mJy} bin exceeds the \SI{244}{mJy} bin by approximately
$2.75\times$, $2.53\times$ and $2.17\times$ at 250, 350 and 500~$\upmu$m respectively.
The rise is present directly in the published tables and is not introduced by the
evaluator or a unit conversion.

This uptick is not evidence for or against \popcosmos{}; it is a feature of the
observational ruler. Valiante discuss larger field-to-field differences at the brightest
fluxes and associate them with cosmic variance. Rare nearby galaxies, small-number
statistics, gravitational lensing, binning effects and the inversion/correction
procedure are all plausible contributors. The available analysis does not isolate their
individual contributions. The three band ratios are consistent with each other within
their uncertainties, so no significant wavelength trend is claimed. This result
motivates caution above approximately \SI{100}{mJy} and shows that validating the
observational inputs was part of the work.

\subsection{Robustness and alternative explanations}
\label{sec:results:robustness}

Five potential concerns are addressed in turn.

\textbf{Could one count survey drive the result?} The block bootstrap over independent
survey families keeps the FSPS offset positive (median $+\SI{0.29}{dex}$; 95\%
bootstrap interval $[+0.16,+0.52]$). This is treated as a robustness interval rather
than a formal confidence statement. The clean three-family evaluator
(Valiante, Oliver, Pearson) gives similar rankings to the full nine-source compilation.
Removing any single survey family does not change the sign of the excess.

\textbf{Could the Wang deblending be responsible for the scatter?} Wang and Jin agree
with each other to 0.12--\SI{0.15}{dex}, while \popcosmos{} scatters by
0.42--\SI{0.51}{dex} against either catalogue. Deblending accounts for 7\% of the
variance at 250~$\upmu$m and 13\% at 500~$\upmu$m; the large majority belongs to the
model.

\textbf{Could SNR selection explain it?} The per-object scatter is stable under SNR
cuts of 3, 5, 7 and $10\sigma$. If noise near the detection limit were responsible,
tighter cuts would shrink the scatter, and they do not.

\textbf{Could AGN dominate the cold-dust result?} The Drew--Casey peak-relation
comparison uses a low-AGN sample ($f_\mathrm{AGN} < 0.1$), which excludes the
hot-peaking branch. The cold, flat-slope result holds for this clean sample. The AGN
issue is discussed separately in Appendix~B and Section~\ref{sec:discussion}.

\textbf{Could lensing dominate?} The central count claim is restricted to the
10--\SI{100}{mJy} range, where lensing corrections are small. The sparse bins above
\SI{100}{mJy}, including the Valiante uptick, are treated as a limitation rather than
the main model constraint.

\section{Discussion}
\label{sec:discussion}

\subsection{What the validation says about \popcosmos{}}
\label{sec:discussion:summary}

The optical and near-infrared population model extrapolates into the far-infrared with
useful total infrared luminosities, but its present far-infrared predictions are not
reliable without an improved dust-emission treatment. Two separate deficiencies are
responsible, and they require different corrections.

The first is a population-level SED-shape problem. The baseline dust emission peaks
near \SI{136}{\um} in the rest frame and is nearly independent of luminosity. This
places too much of a given \LIR{} at SPIRE wavelengths, raising the predicted 250, 350
and 500~$\upmu$m counts above every independent published dataset. Substituting warmer
dust shapes at fixed \LIR{} improves the count agreement, confirming that the shape is
the lever.

The second is an object-level scatter problem. Individual predicted fluxes scatter by
0.44--\SI{0.54}{dex} against both the Wang and Jin catalogues. Changing only
the mean far-infrared template does not reduce this scatter. A single replacement
template therefore cannot solve both problems.

\subsection{Physical interpretation of the cold and uniform FIR shape}
\label{sec:discussion:cold}

Energy balance fixes the integral of dust emission from 8 to \SI{1000}{\um}, but it
does not determine how that energy is distributed over wavelength. This work does not
show that \LIR{} itself is wrong. It shows that the default dust-shape assumptions in
\textsc{fsps} are too restrictive for far-infrared prediction.

The Drew and Casey diagnostic makes this concrete \citep{Drew_2022}. Observed galaxies
follow a clear trend in which more luminous systems have warmer dust that peaks at
shorter wavelengths. Low-AGN \popcosmos{} galaxies do not follow this trend. Their
peak wavelength is nearly flat across four decades of \LIR{}, and it sits roughly
\SI{34}{\um} longer (colder) than the observed relation at
$L_\mathrm{IR} = 10^{12}\,L_\odot$. This missing luminosity--temperature dependence
provides a physical explanation for why the count excess is strongest at 350 and
500~$\upmu$m, where cold dust contributes the most flux.

The controlled template substitutions support the same conclusion from a different
direction. Modified blackbodies, ALESS hybrids and Casey-style models all move the
counts toward agreement when given a warmer shape at fixed \LIR{}
\citep{Casey_2012,Swinbank_2013}. However, several template families score similarly
within the inter-survey scatter, so the counts identify a needed direction (warmer and
luminosity-dependent) but do not uniquely select a dust prescription. The best
chi-square score should not be read as a physical identification of the correct grain
model.

The hypothesis supported by this work is that \popcosmos{} needs a galaxy-dependent
far-infrared shape, likely linked to \LIR{} and potentially to redshift, compactness
or radiation-field intensity, rather than one effectively frozen cold-dust shape.

\subsection{Why the number counts and individual fluxes appear to disagree}
\label{sec:discussion:paradox}

Published counts say the model produces too many bright objects. Matched Wang and Jin
galaxies say the model is typically too faint. Both are true, and the resolution lies
in the shape of the source distribution.

The per-object scatter of approximately \SI{0.45}{dex} moves galaxies both upward and
downward in flux. But because faint galaxies greatly outnumber bright ones, far more
objects scatter upward across a bright threshold than scatter downward across it. This
asymmetric promotion inflates the bright counts even while the median matched galaxy is
underpredicted. At 250~$\upmu$m the scatter effect accounts for the full measured
excess; at 350 and 500~$\upmu$m it accounts for roughly half to two-thirds, with the
cold-dust mean offset contributing the remainder.

The Wang--Jin agreement serves as a control. The two deblending catalogues agree with
each other to 0.12--\SI{0.15}{dex}, far tighter than either agrees with \popcosmos{}
at 0.44--\SI{0.54}{dex}. One deblending pipeline is therefore unlikely to explain the
scatter. Wang and Jin do share the COSMOS field and related prior information, so their
agreement is a cross-check rather than fully independent validation.

The implication is that population summaries and per-object residuals test different
properties of the model. Matching counts alone could hide poor individual predictions,
and matching the median object alone could hide count inflation. Future validation needs
both.

\subsection{How strong are the conclusions?}
\label{sec:discussion:strength}

Several conclusions are robust. The sign of the baseline count excess is stable across
all three independent survey families. The excess grows toward longer SPIRE wavelengths.
Warmer template shapes improve the agreement. Large per-object scatter appears against
both Wang and Jin. These hold under changes to the error floor, the source selection and
the bootstrap resampling.

Other conclusions are more tentative. The exact preferred temperature or template family
depends on the error floor and the flux range used. A formal significance from
chi-square is limited by correlated bins within each survey. The detailed role of AGN
rests on posterior medians of a parameter with a bimodal prior
\citep{Thorp_2025_insights}. Conclusions drawn above \SI{100}{mJy} are weakened by the
small COSMOS area, cosmic variance and the increasing contribution of gravitationally
lensed sources.

The main limitations and their effects are as follows. Only three independent survey
families are available, with correlated bins inside each, so the error floor and block
bootstrap are robustness tools rather than a full likelihood. The COSMOS area of
\SI{1.278}{\deg\squared} gives few very bright sources. Nominal SPIRE wavelengths were
sampled rather than integrating through full filter transmission curves. SEDs use
posterior-median parameters rather than full posterior propagation, which may hide
bimodality in $f_\mathrm{AGN}$ and other parameters. The ALESS template represents
luminous submillimetre galaxies and is not a universal model
\citep{Swinbank_2013}. None of the experiments retrain the population prior jointly
with far-infrared data; they replace or perturb the far-infrared extension after
inference. These limitations weaken model selection more than the basic conclusion that
the baseline far-infrared treatment needs improvement. The practical reliability range
for the count comparison is approximately 10--\SI{100}{mJy}.

\subsection{The Valiante bright-end uptick}
\label{sec:discussion:valiante}

The uptick in the brightest Valiante bins (Section~\ref{sec:results:valiante}) was
checked directly against the published tables and appears in all three bands and all
three GAMA fields. It is not a transcription or unit-conversion error.

This matters because an unusual feature in the observational benchmark can otherwise be
mistaken for a model failure. Small-number statistics, cosmic variance, rare nearby
galaxies, gravitational lensing and the count-inversion procedure are all plausible
contributors, and this project cannot distinguish them \citep{Valiante_2016}. The
correct conclusion is caution above approximately \SI{100}{mJy}, not a claimed new
galaxy population. It illustrates why the sparse bright end should not dominate model
ranking and why published-bin covariance matters.

\subsection{Implications and recommended extension of \popcosmos{}}
\label{sec:discussion:recommendations}

The findings translate into a prioritised set of recommendations for the \popcosmos{}
team.

The most immediate change is to replace the fixed far-infrared bump with a
galaxy-dependent physical dust-emission family. Variable Draine and Li parameters, such
as radiation-field intensity and warm-dust fraction \citep{Draine_2007}, or alternative
dust modules available through CIGALE, would allow the far-infrared shape to respond to
galaxy properties. The Casey and ALESS templates used here are directional diagnostics
and useful brackets, not recommended endpoint models.

At the population level, a prior or conditional relation between far-infrared shape and
properties such as \LIR{}, redshift and SFR should be introduced, guided initially by
the observed $L_\mathrm{IR}$--$\lambda_\mathrm{peak}$ relation. This relation should
include scatter rather than assigning one deterministic template to every galaxy at a
given luminosity.

For validation, a calibrate-then-validate loop would strengthen any future claim.
Calibrate on one set of independent survey fields and evaluate on held-out fields.
Score differential counts by band and flux regime, and separately report per-object bias
and scatter against deblended catalogues. Redshift-sliced counts or stacking
constraints from \citet{B_thermin_2012} would test whether a correction is being applied
to the right population.

Technical improvements include integrating through the true SPIRE bandpasses,
propagating posterior samples for the high-AGN and extreme-SFR populations, and
modelling survey selection and covariance explicitly where possible.

This work extends \popcosmos{} validation beyond its training wavelengths, identifies
where and why the extrapolation fails, and supplies a reproducible observed-space
evaluation framework for testing the next physical dust prescription.

\section{Conclusions}
\label{sec:conclusions}

This project tested whether \popcosmos{}, a generative galaxy population model
constrained at optical and near-infrared wavelengths, makes reliable predictions in the
\emph{Herschel}/SPIRE bands at 250, 350 and \SI{500}{\um}. The validation used
published differential number counts and matched COSMOS catalogues as directly observed
benchmarks, without relying on derived physical quantities. The model works well where
it is constrained, but its current far-infrared extension does not reliably reproduce
either the population counts or the individual galaxy fluxes.

Nine published count products from six papers were placed in common units and grouped
into three independent survey families, with overlapping products retained as robustness
checks. The baseline \popcosmos{}/FSPS prediction overpredicts the counts in all three
bands. Over the 30--\SI{100}{mJy} range where all three families overlap, the excess is
approximately $1.7\times$ at 250 and 350~$\upmu$m and $3.6\times$ at 500~$\upmu$m. As a
control, IRAC Ch1 and Ch2 residuals are only a few hundredths of a magnitude, confirming
that the mismatch is wavelength-specific rather than a broken catalogue or flux pipeline.

The physical diagnosis is that the baseline far-infrared dust bump is too cold and too
uniform. At $L_\mathrm{IR} = 10^{12}\,L_\odot$, model galaxies peak near
\SI{126}{\um} rather than the observed \SI{92}{\um}, and the model shows essentially no
luminosity dependence (slope $\eta = +0.005$ vs the observed $-0.09$). Preserving each
galaxy's total \LIR{} while substituting warmer dust shapes (modified blackbodies,
ALESS/FSPS hybrids, Casey-style models) improves the counts. Their similar performance
identifies the required direction (warmer and luminosity-dependent) but does not
uniquely select a physical dust model.

A separate problem emerged from the per-object comparison. Model fluxes scatter by
0.44--\SI{0.54}{dex} against the Wang catalogue, while Wang and Jin agree with each
other to only 0.12--\SI{0.15}{dex}. The median matched galaxy is underpredicted, yet
the model still overproduces bright counts, because the steep source distribution and
large scatter promote many more faint objects upward than bright objects downward.
Changing the mean dust template improves the counts but does not remove this excessive
object-to-object scatter. These are two deficiencies, not one.

The excess keeps the same sign across all three independent survey families, under
changes to the error floor and source selection, and the per-object scatter persists
under SNR cuts from 3 to $10\sigma$. The block bootstrap gives a median baseline offset
of $+\SI{0.29}{dex}$ with an interval of $[+0.16, +0.52]$, treated as a robustness
check rather than a formal exclusion. The main limitations are correlated bins within
each survey, only three independent families, nominal-wavelength rather than
full-bandpass fluxes, posterior medians rather than full posterior propagation, and weak
bright-end statistics in a COSMOS-sized area. The Valiante bright-end uptick was
verified as a feature of the published counts and motivates caution above approximately
\SI{100}{mJy}; its origin was not resolved.

The primary recommendation for the \popcosmos{} team is to preserve the learned
population model and energy balance, but introduce a galaxy-dependent far-infrared dust
family linked to \LIR{} and potentially redshift and SFR, with intrinsic scatter.
Variable Draine and Li parameters \citep{Draine_2007} or alternative dust modules
available through CIGALE are natural starting points. Future validation should calibrate
on one set of survey fields and test on held-out fields, reporting both differential
count performance and per-object flux scatter. The general lesson is that a generative
population model can produce physically coherent parameters while still failing in
observables beyond its constraining wavelength range. Out-of-sample validation needs
both population distributions and individual-object tests.

\section*{Acknowledgements}
% TODO: supervisor, Boris (FSPS SED and MIR emulator outputs), the pop-cosmos team.

% ===================== AI USE DECLARATION =====================
% Imperial requires this to be accurate and specific. EDIT to reflect exactly what
% you used and how. A vague declaration is worse than none.
\FloatBarrier
\section*{AI Use Declaration}
Generative AI tools were used in the preparation of this report in the following
specific ways: <<EDIT --- e.g.\ drafting and editing of prose subsequently rewritten by
the author; assistance in interpreting technical details of cited papers; independent
re-derivation and cross-checking of numerical results reported in
Section~\ref{sec:results}; LaTeX formatting and figure layout>>. All scientific
conclusions, the choice of analysis methods, and the interpretation of results are the
author's own. Numerical results were produced by code written for this project
<<add repository reference>> and independently verified.

% ===================== BIBLIOGRAPHY =====================
\bibliographystyle{apalike}
\bibliography{references}

\end{document}
